\documentclass[11pt,a4paper]{article}
\usepackage[margin=25mm,headheight=14pt]{geometry}
\usepackage{amsmath,amssymb,amsthm}
\usepackage{fontspec}
\usepackage{unicode-math}
\ifChinese
 \IfFontExistsTF{Noto Serif CJK SC}{\setmainfont{Noto Serif CJK SC}[ItalicFont={Noto Serif CJK SC},BoldFont={Noto Serif CJK SC}]}{\setmainfont{SimSun}[ItalicFont={SimSun},BoldFont={SimHei}]}
 \IfFontExistsTF{Noto Sans CJK SC}{\setsansfont{Noto Sans CJK SC}}{\setsansfont{Microsoft YaHei}}
 \IfFontExistsTF{Noto Sans Mono CJK SC}{\setmonofont{Noto Sans Mono CJK SC}}{\setmonofont{SimSun}[BoldFont={SimHei}]}
 \XeTeXlinebreaklocale "zh"
 \XeTeXlinebreakskip=0pt plus 1pt
 \newcommand{\bi}[2]{#2}
\else
 \IfFontExistsTF{TeX Gyre Termes}{\setmainfont{TeX Gyre Termes}}{\setmainfont{Times New Roman}}
 \IfFontExistsTF{TeX Gyre Heros}{\setsansfont{TeX Gyre Heros}}{\setsansfont{Arial}}
 \IfFontExistsTF{DejaVu Sans Mono}{\setmonofont{DejaVu Sans Mono}[Scale=0.9]}{\setmonofont{Consolas}[Scale=0.9]}
 \newcommand{\bi}[2]{#1}
\fi
\setmathfont{Latin Modern Math}[BoldFont={Latin Modern Math}]
\usepackage{booktabs,longtable,array}
\IfFileExists{xurl.sty}{\usepackage{xurl}}{\usepackage{url}}
\usepackage{xcolor}
\usepackage[colorlinks=true,linkcolor=blue!50!black,citecolor=blue!50!black,urlcolor=blue!50!black]{hyperref}
\IfFileExists{microtype.sty}{\usepackage{microtype}}{}
\newtheorem{theorem}{\bi{Theorem}{定理}}[section]
\newtheorem{lemma}[theorem]{\bi{Lemma}{引理}}
\newtheorem{proposition}[theorem]{\bi{Proposition}{命题}}
\newtheorem{corollary}[theorem]{\bi{Corollary}{推论}}
\theoremstyle{definition}
\newtheorem{definition}[theorem]{\bi{Definition}{定义}}
\theoremstyle{remark}
\newtheorem{remark}[theorem]{\bi{Remark}{注}}
\renewcommand{\proofname}{\bi{Proof}{证明}}
\renewcommand{\abstractname}{\bi{Abstract}{摘要}}
\renewcommand{\contentsname}{\bi{Contents}{目录}}
\renewcommand{\refname}{\bi{References}{参考文献}}
\newcommand{\Tr}{\operatorname{Tr}}
\newcommand{\Sym}{\operatorname{Sym}}
\newcommand{\CM}{\mathrm{COM}}
\newcommand{\ran}{\operatorname{ran}}
\newcommand{\Spec}{\operatorname{Spec}}
\newcommand{\dd}{\,\mathrm d}
\newcommand{\ket}[1]{|#1\rangle}
\newcommand{\bra}[1]{\langle#1|}
\newcommand{\norm}[1]{\lVert#1\rVert}
\newcommand{\abs}[1]{\lvert#1\rvert}
\newcommand{\inner}[2]{\langle#1,#2\rangle}
\allowdisplaybreaks[2]
\emergencystretch=2em
\setlength{\parskip}{3pt plus 1pt}
\setlength{\parindent}{1.2em}
\numberwithin{equation}{section}
\setcounter{tocdepth}{2}
\hypersetup{pdftitle={Uniform Laughlin gaps on large flat tori from finite Fock-space certificates},pdfauthor={prepared by Xiang-Jian Hou with advice from Ying-Hai Wu},pdfsubject={Candidate proof for independent review: bosonic V0 and fermionic V1 on a magnetic torus}}
